\section{Conclusions}
\label{sec:conclusions}

\subsection{Established}
From retained generating circles: Angular Shortfall, membrane body, layered
particle, residual-stress dipole, pressure length, elastic moduli ratios, optical
modes and dual dissipation, contact rule $N_c^*=3$, helical parameters, and the
necessity of cyclicity.

\subsection{Status of main claims}
\begin{center}
\begin{tabular}{ll}
\toprule
Claim & Status \\
\midrule
Angular Shortfall as geometric necessity & Derived \\
Layered particle from four-sphere overlap & Derived \\
Residual-stress dipole $P/p$ & Derived \\
Pressure length $\ell=s/\sqrt{10}$; regimes $\mathrm{Pl},\mathrm{pL},\mathrm{PL},\mathrm{pl}$ & Derived \\
Elastic moduli ratios $Y,K,G$ & Derived (ratios) \\
Five optical modes & Conceptual (geometry-seated) \\
Dual dissipation channels & Conceptual \\
Scale closure with $\hbar c$, $s\approx 0.84\,\mathrm{fm}$ & Assumed / interpretive \\
$N_c^*=3$ pathways & Derived \\
Cyclicity forced by shortfall & Derived \\
EM field from membrane leakage & Motivated hypothesis \\
Quantitative cosmology & Open \\
\bottomrule
\end{tabular}
\end{center}

\subsection{Open questions}
\begin{enumerate}
\item Precise geometric prefactor $C$ in $|\sigma_{\mathrm{res}}|=C\,\hbar c/s^2$.
\item Full 3D ray-trace with shortfall-controlled deflection.
\item Detailed laboratory mapping of $P/p$ and $N/n$.
\item Whether $s$ is exactly the proton radius or a geometric multiple thereof.
\item Quantitative cosmological confrontation.
\end{enumerate}

\subsection{Final remark}
The Angular Shortfall is small, permanent, and present at every scale. From that
single fact the layered particle, residual stress, optical organisation, helical
structure and the necessity of cycles all follow. The universe described here is
never at rest, not because of an external driver, but because perfect closure was
never available.

\appendix
